
\documentclass{article}
\usepackage{colortbl}
\usepackage{makecell}
\usepackage{multirow}
\usepackage{supertabular}

\begin{document}

\newcounter{utterance}

\centering \large Interaction Transcript for game `function\_detective', experiment `list\_passive\_examples', episode 67 with gpt{-}5.4.
\vspace{24pt}

{ \footnotesize  \setcounter{utterance}{1}
\setlength{\tabcolsep}{0pt}
\begin{supertabular}{c@{$\;$}|p{.15\linewidth}@{}p{.15\linewidth}p{.15\linewidth}p{.15\linewidth}p{.15\linewidth}p{.15\linewidth}}
    \# & \multicolumn{2}{c}{Player} && \multicolumn{2}{c}{Game Master} \\
    \hline

    \theutterance \stepcounter{utterance}  
    & & & \multicolumn{4}{p{0.6\linewidth}}{
        \cellcolor[rgb]{0.9,0.9,0.9}{
            \makecell[{{p{\linewidth}}}]{
                \texttt{\tiny{[P1$\langle$GM]}}
                \texttt{I have a hidden Python function for you to figure out.} \\
\\ 
\texttt{Your job is to infer what it does and then submit a correct Python implementation.} \\
\\ 
\texttt{Function signature:} \\
\texttt{(items: List{[}int{]}) {-}> int} \\
\\ 
\texttt{You have 15 turns total.} \\
\\ 
\texttt{I control which valid input{-}output examples you see. You can request them one at a time.} \\
\\ 
\texttt{You must make your best SOLVE attempt no later than the second{-}to{-}last turn.} \\
\texttt{Do not wait until the very end.} \\
\\ 
\texttt{You can use these actions:} \\
\texttt{{-} JSON format described at the end of the prompt} \\
\texttt{{-} SOLVE: ```python ... ```} \\
\\ 
\texttt{Do not add import statements in your guess.} \\
\texttt{Anything you might need is already available, so just write the function body in hypothesis and return the answer.} \\
\\ 
\texttt{Preloaded examples:} \\
\texttt{1. ({[}{-}1, 0, 1{]}, 3)} \\
\\ 
\texttt{Scoring:} \\
\texttt{{-} correctness is required for any non{-}zero score} \\
\texttt{{-} in interactive modes, solving with fewer information requests is better} \\
\\ 
\texttt{Input formatting reminders:} \\
\texttt{{-} strings must be quoted} \\
\texttt{{-} booleans must be True/False} \\
\texttt{{-} use commas between multiple values} \\
\texttt{{-} do not write any extra action tags in the same message} \\
\\ 
\texttt{Format example:} \\
\texttt{You write:} \\
\texttt{\{'current\_mode': <probing/confirming/don't know; choose whichever best describes your current state in the task>, 'input\_rationale': <brief description of why you chose this hypothesis>, 'hypothesis': <either the exact string "unknown" if you have no guess yet OR your current best guess as runnable Python matching the given signature, with a one{-}line docstring stating what it does. Never write placeholder code, if you cannot commit to a guess, use "unknown". Match the given signature.} \\
\texttt{e.g.} \\
\texttt{def f(x):} \\
\texttt{	"""Returns the absolute value of x."""} \\
\texttt{	return x if x >= 0 else {-}x} \\
\\ 
\texttt{Use "unknown" if you have no guess.>, 'confidence\_score': <score between 1{-}100. 100 means highest confidence in your hypothesis>\}} \\
\\ 
\texttt{Output only one valid json per turn.} \\
            }
        }
    }
     \\ \\

    \theutterance \stepcounter{utterance}  
    & \multicolumn{4}{p{0.6\linewidth}}{
        \cellcolor[rgb]{0.561,0.737,0.561}{
            \makecell[{{p{\linewidth}}}]{
                \texttt{\tiny{[P1$\rangle$GM]}}
                \texttt{\{"current\_mode":"probing","input\_rationale":"Need a more discriminating example; many simple functions fit {[}{-}1,0,1{]} {-}> 3 (e.g. len, sum of abs, range width + 1, distinct count, etc.). Please provide another valid input{-}output example, ideally with duplicates or all{-}positive values.","hypothesis":"unknown","confidence\_score":5\}} \\
            }
        }
    }
    & & \\ \\

    \theutterance \stepcounter{utterance}  
    & & & \multicolumn{4}{p{0.6\linewidth}}{
        \cellcolor[rgb]{0.9,0.9,0.9}{
            \makecell[{{p{\linewidth}}}]{
                \texttt{\tiny{[P1$\langle$GM]}}
                \texttt{OUTPUT: ({[}5, {-}5, 5{]}, 5)} \\
            }
        }
    }
     \\ \\

    \theutterance \stepcounter{utterance}  
    & \multicolumn{4}{p{0.6\linewidth}}{
        \cellcolor[rgb]{0.561,0.737,0.561}{
            \makecell[{{p{\linewidth}}}]{
                \texttt{\tiny{[P1$\rangle$GM]}}
                \texttt{\{"current\_mode":"probing","input\_rationale":"Current examples still fit many possibilities. A case with repeated zeros or a short two{-}element list would sharply narrow it. Please provide another valid input{-}output example, ideally one of {[}0,0{]}, {[}2,2{]}, {[}2,{-}2{]}, or a similar small case.","hypothesis":"unknown","confidence\_score":9\}} \\
            }
        }
    }
    & & \\ \\

    \theutterance \stepcounter{utterance}  
    & & & \multicolumn{4}{p{0.6\linewidth}}{
        \cellcolor[rgb]{0.9,0.9,0.9}{
            \makecell[{{p{\linewidth}}}]{
                \texttt{\tiny{[P1$\langle$GM]}}
                \texttt{OUTPUT: ({[}1, 1{]}, 2)} \\
            }
        }
    }
     \\ \\

    \theutterance \stepcounter{utterance}  
    & \multicolumn{4}{p{0.6\linewidth}}{
        \cellcolor[rgb]{0.561,0.737,0.561}{
            \makecell[{{p{\linewidth}}}]{
                \texttt{\tiny{[P1$\rangle$GM]}}
                \texttt{\{"current\_mode":"probing","input\_rationale":"A tiny edge case with zeros or opposite signs should eliminate many remaining hypotheses. Please provide one more valid example, ideally with two elements such as {[}0,0{]}, {[}2,{-}2{]}, or {[}2,3{]}.","hypothesis":"unknown","confidence\_score":12\}} \\
            }
        }
    }
    & & \\ \\

    \theutterance \stepcounter{utterance}  
    & & & \multicolumn{4}{p{0.6\linewidth}}{
        \cellcolor[rgb]{0.9,0.9,0.9}{
            \makecell[{{p{\linewidth}}}]{
                \texttt{\tiny{[P1$\langle$GM]}}
                \texttt{OUTPUT: ({[}0{]}, 0)} \\
            }
        }
    }
     \\ \\

    \theutterance \stepcounter{utterance}  
    & \multicolumn{4}{p{0.6\linewidth}}{
        \cellcolor[rgb]{0.561,0.737,0.561}{
            \makecell[{{p{\linewidth}}}]{
                \texttt{\tiny{[P1$\rangle$GM]}}
                \texttt{\{"current\_mode":"confirming","input\_rationale":"I have a concrete hypothesis: the function may return the absolute value of the sum of running prefix minima. A short two{-}element example like {[}2,3{]} or {[}2,{-}2{]} will confirm or refute it cleanly. Please provide one more valid input{-}output example, ideally one of those.","hypothesis":"def f(items):\textbackslash{}n    \textbackslash{}"\textbackslash{}"\textbackslash{}"Returns abs(sum of the minimum value seen so far at each position).\textbackslash{}"\textbackslash{}"\textbackslash{}"\textbackslash{}n    if not items:\textbackslash{}n        return 0\textbackslash{}n    m = items{[}0{]}\textbackslash{}n    s = 0\textbackslash{}n    for x in items:\textbackslash{}n        if x < m:\textbackslash{}n            m = x\textbackslash{}n        s += m\textbackslash{}n    return s if s >= 0 else {-}s","confidence\_score":54\}} \\
            }
        }
    }
    & & \\ \\

    \theutterance \stepcounter{utterance}  
    & & & \multicolumn{4}{p{0.6\linewidth}}{
        \cellcolor[rgb]{0.9,0.9,0.9}{
            \makecell[{{p{\linewidth}}}]{
                \texttt{\tiny{[P1$\langle$GM]}}
                \texttt{OUTPUT: ({[}7{]}, 1)} \\
            }
        }
    }
     \\ \\

    \theutterance \stepcounter{utterance}  
    & \multicolumn{4}{p{0.6\linewidth}}{
        \cellcolor[rgb]{0.561,0.737,0.561}{
            \makecell[{{p{\linewidth}}}]{
                \texttt{\tiny{[P1$\rangle$GM]}}
                \texttt{\{"current\_mode":"probing","input\_rationale":"The singleton cases suggest behavior depends strongly on sign/value class. A negative singleton or a simple two{-}element opposite{-}sign case would be highly informative. Please provide one more valid example, ideally either ({[}{-}3{]}, ?), ({[}5{]}, ?), or ({[}2, {-}2{]}, ?).","hypothesis":"unknown","confidence\_score":14\}} \\
            }
        }
    }
    & & \\ \\

    \theutterance \stepcounter{utterance}  
    & & & \multicolumn{4}{p{0.6\linewidth}}{
        \cellcolor[rgb]{0.9,0.9,0.9}{
            \makecell[{{p{\linewidth}}}]{
                \texttt{\tiny{[P1$\langle$GM]}}
                \texttt{OUTPUT: ({[}2, 1{]}, 3)} \\
            }
        }
    }
     \\ \\

    \theutterance \stepcounter{utterance}  
    & \multicolumn{4}{p{0.6\linewidth}}{
        \cellcolor[rgb]{0.561,0.737,0.561}{
            \makecell[{{p{\linewidth}}}]{
                \texttt{\tiny{[P1$\rangle$GM]}}
                \texttt{\{"current\_mode":"probing","input\_rationale":"The current data still supports several very different rules. A negative singleton or a mixed{-}sign two{-}element case would separate them quickly. Please provide one more example, ideally either ({[}{-}3{]}, ?), ({[}2, {-}2{]}, ?), or ({[}3, 3{]}, ?).","hypothesis":"unknown","confidence\_score":16\}} \\
            }
        }
    }
    & & \\ \\

    \theutterance \stepcounter{utterance}  
    & & & \multicolumn{4}{p{0.6\linewidth}}{
        \cellcolor[rgb]{0.9,0.9,0.9}{
            \makecell[{{p{\linewidth}}}]{
                \texttt{\tiny{[P1$\langle$GM]}}
                \texttt{OUTPUT: ({[}{]}, 0)} \\
            }
        }
    }
     \\ \\

    \theutterance \stepcounter{utterance}  
    & \multicolumn{4}{p{0.6\linewidth}}{
        \cellcolor[rgb]{0.561,0.737,0.561}{
            \makecell[{{p{\linewidth}}}]{
                \texttt{\tiny{[P1$\rangle$GM]}}
                \texttt{\{"current\_mode":"confirming","input\_rationale":"A strong current hypothesis is: remove zeros, then return the number of nonzero elements plus the number of times the value changes between consecutive remaining elements. To test whether equality across a zero is treated as the same run, please provide one more example such as ({[}1, 0, 1{]}, ?), ({[}2, 2{]}, ?), or ({[}2, {-}2{]}, ?).","hypothesis":"def f(items):\textbackslash{}n    \textbackslash{}"\textbackslash{}"\textbackslash{}"Removes zeros and returns count(nonzero items) plus count of value changes between consecutive remaining items.\textbackslash{}"\textbackslash{}"\textbackslash{}"\textbackslash{}n    prev = None\textbackslash{}n    seen = False\textbackslash{}n    ans = 0\textbackslash{}n    for x in items:\textbackslash{}n        if x == 0:\textbackslash{}n            continue\textbackslash{}n        ans += 1\textbackslash{}n        if seen and x != prev:\textbackslash{}n            ans += 1\textbackslash{}n        prev = x\textbackslash{}n        seen = True\textbackslash{}n    return ans","confidence\_score":62\}} \\
            }
        }
    }
    & & \\ \\

    \theutterance \stepcounter{utterance}  
    & & & \multicolumn{4}{p{0.6\linewidth}}{
        \cellcolor[rgb]{0.9,0.9,0.9}{
            \makecell[{{p{\linewidth}}}]{
                \texttt{\tiny{[P1$\langle$GM]}}
                \texttt{OUTPUT: ({[}3, 2, 1{]}, 3)} \\
            }
        }
    }
     \\ \\

    \theutterance \stepcounter{utterance}  
    & \multicolumn{4}{p{0.6\linewidth}}{
        \cellcolor[rgb]{0.561,0.737,0.561}{
            \makecell[{{p{\linewidth}}}]{
                \texttt{\tiny{[P1$\rangle$GM]}}
                \texttt{\{"current\_mode":"confirming","input\_rationale":"Most evidence supports a rule based on counting nonzero entries plus sign flips across consecutive nonzero values, but {[}2,1{]} is an outlier. A short alternating{-}sign example would strongly confirm or refute that family. Please provide one more example, ideally ({[}1, {-}1, 1{]}, ?).","hypothesis":"def f(items):\textbackslash{}n    \textbackslash{}"\textbackslash{}"\textbackslash{}"Returns the number of nonzero items plus the number of sign changes between consecutive nonzero items.\textbackslash{}"\textbackslash{}"\textbackslash{}"\textbackslash{}n    prev = 0\textbackslash{}n    ans = 0\textbackslash{}n    for x in items:\textbackslash{}n        if x == 0:\textbackslash{}n            continue\textbackslash{}n        ans += 1\textbackslash{}n        if prev and ((x > 0) != (prev > 0)):\textbackslash{}n            ans += 1\textbackslash{}n        prev = x\textbackslash{}n    return ans","confidence\_score":41\}} \\
            }
        }
    }
    & & \\ \\

    \theutterance \stepcounter{utterance}  
    & & & \multicolumn{4}{p{0.6\linewidth}}{
        \cellcolor[rgb]{0.9,0.9,0.9}{
            \makecell[{{p{\linewidth}}}]{
                \texttt{\tiny{[P1$\langle$GM]}}
                \texttt{OUTPUT: ({[}1, 2, 3{]}, 3)} \\
            }
        }
    }
     \\ \\

    \theutterance \stepcounter{utterance}  
    & \multicolumn{4}{p{0.6\linewidth}}{
        \cellcolor[rgb]{0.561,0.737,0.561}{
            \makecell[{{p{\linewidth}}}]{
                \texttt{\tiny{[P1$\rangle$GM]}}
                \texttt{\{"current\_mode":"confirming","input\_rationale":"The strongest consistent pattern is: ignore zeros, count each nonzero element, and add 1 for each sign change between consecutive nonzero elements. All examples fit except {[}2,1{]}, so I need one direct same{-}sign control to decide whether that case is special or anomalous. Please provide one more example, ideally ({[}2, 2{]}, ?) or ({[}2, {-}2{]}, ?).","hypothesis":"def f(items):\textbackslash{}n    \textbackslash{}"\textbackslash{}"\textbackslash{}"Returns the number of nonzero items plus the number of sign changes between consecutive nonzero items.\textbackslash{}"\textbackslash{}"\textbackslash{}"\textbackslash{}n    prev = 0\textbackslash{}n    ans = 0\textbackslash{}n    for x in items:\textbackslash{}n        if x == 0:\textbackslash{}n            continue\textbackslash{}n        ans += 1\textbackslash{}n        if prev != 0 and ((x > 0) != (prev > 0)):\textbackslash{}n            ans += 1\textbackslash{}n        prev = x\textbackslash{}n    return ans","confidence\_score":58\}} \\
            }
        }
    }
    & & \\ \\

    \theutterance \stepcounter{utterance}  
    & & & \multicolumn{4}{p{0.6\linewidth}}{
        \cellcolor[rgb]{0.9,0.9,0.9}{
            \makecell[{{p{\linewidth}}}]{
                \texttt{\tiny{[P1$\langle$GM]}}
                \texttt{OUTPUT: ({[}1, 10, {-}10{]}, 4)} \\
            }
        }
    }
     \\ \\

    \theutterance \stepcounter{utterance}  
    & \multicolumn{4}{p{0.6\linewidth}}{
        \cellcolor[rgb]{0.561,0.737,0.561}{
            \makecell[{{p{\linewidth}}}]{
                \texttt{\tiny{[P1$\rangle$GM]}}
                \texttt{\{"current\_mode":"probing","input\_rationale":"There is a near{-}consistent sign{-}change hypothesis, but {[}2, 1{]} is a crucial exception. A same{-}sign control will tell me whether that exception is real structure or noise. Please provide one more example, ideally ({[}2, 2{]}, ?) or ({[}2, {-}2{]}, ?).","hypothesis":"unknown","confidence\_score":22\}} \\
            }
        }
    }
    & & \\ \\

    \theutterance \stepcounter{utterance}  
    & & & \multicolumn{4}{p{0.6\linewidth}}{
        \cellcolor[rgb]{0.9,0.9,0.9}{
            \makecell[{{p{\linewidth}}}]{
                \texttt{\tiny{[P1$\langle$GM]}}
                \texttt{OUTPUT: ({[}1{]}, 1)} \\
            }
        }
    }
     \\ \\

    \theutterance \stepcounter{utterance}  
    & \multicolumn{4}{p{0.6\linewidth}}{
        \cellcolor[rgb]{0.561,0.737,0.561}{
            \makecell[{{p{\linewidth}}}]{
                \texttt{\tiny{[P1$\rangle$GM]}}
                \texttt{\{"current\_mode":"probing","input\_rationale":"One crucial ambiguity remains: whether any decreasing positive pair gets the extra count seen in {[}2, 1{]}. A direct pair test will let me finalize the rule. Please provide one more example, ideally ({[}3, 2{]}, ?).","hypothesis":"unknown","confidence\_score":24\}} \\
            }
        }
    }
    & & \\ \\

    \theutterance \stepcounter{utterance}  
    & & & \multicolumn{4}{p{0.6\linewidth}}{
        \cellcolor[rgb]{0.9,0.9,0.9}{
            \makecell[{{p{\linewidth}}}]{
                \texttt{\tiny{[P1$\langle$GM]}}
                \texttt{OUTPUT: ({[}{-}2, {-}40{]}, 3)} \\
            }
        }
    }
     \\ \\

    \theutterance \stepcounter{utterance}  
    & \multicolumn{4}{p{0.6\linewidth}}{
        \cellcolor[rgb]{0.561,0.737,0.561}{
            \makecell[{{p{\linewidth}}}]{
                \texttt{\tiny{[P1$\rangle$GM]}}
                \texttt{\{"current\_mode":"probing","input\_rationale":"The critical ambiguity is how two unequal same{-}sign values are handled. A single increasing positive pair would distinguish several plausible rules immediately. Please provide one more example, ideally ({[}2, 3{]}, ?).","hypothesis":"unknown","confidence\_score":19\}} \\
            }
        }
    }
    & & \\ \\

    \theutterance \stepcounter{utterance}  
    & & & \multicolumn{4}{p{0.6\linewidth}}{
        \cellcolor[rgb]{0.9,0.9,0.9}{
            \makecell[{{p{\linewidth}}}]{
                \texttt{\tiny{[P1$\langle$GM]}}
                \texttt{OUTPUT: ({[}{-}10{]}, 2)} \\
            }
        }
    }
     \\ \\

    \theutterance \stepcounter{utterance}  
    & \multicolumn{4}{p{0.6\linewidth}}{
        \cellcolor[rgb]{0.561,0.737,0.561}{
            \makecell[{{p{\linewidth}}}]{
                \texttt{\tiny{[P1$\rangle$GM]}}
                \texttt{SOLVE: ```python} \\
\texttt{def f(items):} \\
\texttt{    """Returns a compact digit{-}count based score derived from the first value and consecutive differences."""} \\
\texttt{    if not items:} \\
\texttt{        return 0} \\
\texttt{    if all(x == 0 for x in items):} \\
\texttt{        return 0} \\
\\ 
\texttt{    def digits(x):} \\
\texttt{        x = {-}x if x < 0 else x} \\
\texttt{        if x == 0:} \\
\texttt{            return 1} \\
\texttt{        c = 0} \\
\texttt{        while x:} \\
\texttt{            c += 1} \\
\texttt{            x //= 10} \\
\texttt{        return c} \\
\\ 
\texttt{    ans = 0 if items{[}0{]} == 0 else digits(items{[}0{]})} \\
\texttt{    for i in range(1, len(items)):} \\
\texttt{        ans += digits(items{[}i{]} {-} items{[}i {-} 1{]})} \\
\\ 
\texttt{    if len(items) >= 2:} \\
\texttt{        d = items{[}1{]} {-} items{[}0{]}} \\
\texttt{        if d in (1, {-}1):} \\
\texttt{            ok = True} \\
\texttt{            for i in range(2, len(items)):} \\
\texttt{                if items{[}i{]} {-} items{[}i {-} 1{]} != d:} \\
\texttt{                    ok = False} \\
\texttt{                    break} \\
\texttt{            if ok:} \\
\texttt{                same\_side = True} \\
\texttt{                for x in items:} \\
\texttt{                    if x == 0 or (x > 0) != (items{[}0{]} > 0):} \\
\texttt{                        same\_side = False} \\
\texttt{                        break} \\
\texttt{                if same\_side:} \\
\texttt{                    ans += 1} \\
\\ 
\texttt{    return ans} \\
\texttt{```} \\
            }
        }
    }
    & & \\ \\

    \theutterance \stepcounter{utterance}  
    & & & \multicolumn{2}{p{0.3\linewidth}}{
        \cellcolor[rgb]{0.9,0.9,0.9}{
            \makecell[{{p{\linewidth}}}]{
                \texttt{\tiny{[GM$|$GM]}}
                \texttt{Game Verdict: LOSS} \\
            }
        }
    }
    & & \\ \\

    \theutterance \stepcounter{utterance}  
    & & & \multicolumn{2}{p{0.3\linewidth}}{
        \cellcolor[rgb]{0.9,0.9,0.9}{
            \makecell[{{p{\linewidth}}}]{
                \texttt{\tiny{[GM$|$GM]}}
                \texttt{Incorrect.} \\
\texttt{Failed Cases:} \\
\texttt{Input: {[}{[}3, 2, 1{]}{]} | Expected: 3 | Got: 4} \\
\texttt{Input: {[}{[}1, 2, 3{]}{]} | Expected: 3 | Got: 4} \\
\texttt{Input: {[}{[}{-}22, 1, {-}10, {-}2{]}{]} | Expected: 6 | Got: 7} \\
            }
        }
    }
    & & \\ \\

    \theutterance \stepcounter{utterance}  
    & & & \multicolumn{2}{p{0.3\linewidth}}{
        \cellcolor[rgb]{0.9,0.9,0.9}{
            \makecell[{{p{\linewidth}}}]{
                \texttt{\tiny{[GM$|$GM]}}
                \texttt{def sum\_mod\_three(items: List{[}int{]}) {-}> int:} \\
\texttt{    return sum(x \% 3 for x in items)} \\
            }
        }
    }
    & & \\ \\

    \theutterance \stepcounter{utterance}  
    & & & \multicolumn{2}{p{0.3\linewidth}}{
        \cellcolor[rgb]{0.9,0.9,0.9}{
            \makecell[{{p{\linewidth}}}]{
                \texttt{\tiny{[GM$|$GM]}}
                \texttt{def f(items):} \\
\texttt{    """Returns a compact digit{-}count based score derived from the first value and consecutive differences."""} \\
\texttt{    if not items:} \\
\texttt{        return 0} \\
\texttt{    if all(x == 0 for x in items):} \\
\texttt{        return 0} \\
\\ 
\texttt{    def digits(x):} \\
\texttt{        x = {-}x if x < 0 else x} \\
\texttt{        if x == 0:} \\
\texttt{            return 1} \\
\texttt{        c = 0} \\
\texttt{        while x:} \\
\texttt{            c += 1} \\
\texttt{            x //= 10} \\
\texttt{        return c} \\
\\ 
\texttt{    ans = 0 if items{[}0{]} == 0 else digits(items{[}0{]})} \\
\texttt{    for i in range(1, len(items)):} \\
\texttt{        ans += digits(items{[}i{]} {-} items{[}i {-} 1{]})} \\
\\ 
\texttt{    if len(items) >= 2:} \\
\texttt{        d = items{[}1{]} {-} items{[}0{]}} \\
\texttt{        if d in (1, {-}1):} \\
\texttt{            ok = True} \\
\texttt{            for i in range(2, len(items)):} \\
\texttt{                if items{[}i{]} {-} items{[}i {-} 1{]} != d:} \\
\texttt{                    ok = False} \\
\texttt{                    break} \\
\texttt{            if ok:} \\
\texttt{                same\_side = True} \\
\texttt{                for x in items:} \\
\texttt{                    if x == 0 or (x > 0) != (items{[}0{]} > 0):} \\
\texttt{                        same\_side = False} \\
\texttt{                        break} \\
\texttt{                if same\_side:} \\
\texttt{                    ans += 1} \\
\\ 
\texttt{    return ans} \\
            }
        }
    }
    & & \\ \\

    \theutterance \stepcounter{utterance}  
    & & & \multicolumn{2}{p{0.3\linewidth}}{
        \cellcolor[rgb]{0.9,0.9,0.9}{
            \makecell[{{p{\linewidth}}}]{
                \texttt{\tiny{[GM$|$GM]}}
                \texttt{Mode: passive\_examples} \\
\texttt{Accuracy: 0.300} \\
\texttt{Information requests: 12/15} \\
\texttt{Revealed passive examples: 13} \\
\texttt{Efficiency: 0.267} \\
\texttt{Slack: 3} \\
\texttt{Observed{-}pair consistency: 0.846} \\
\texttt{Parse errors: 0} \\
\texttt{Runtime errors: 0} \\
            }
        }
    }
    & & \\ \\

\end{supertabular}
}

\end{document}
